\documentclass[10pt, letterpaper]{article}

% Packages:
\usepackage[
    ignoreheadfoot, % set margins without considering header and footer
    top=0.5 cm, % seperation between body and page edge from the top
    bottom=0.5 cm, % seperation between body and page edge from the bottom
    left=0.75 cm, % seperation between body and page edge from the left
    right=0.75 cm, % seperation between body and page edge from the right
    footskip=1.0 cm, % seperation between body and footer
    % showframe % for debugging
]{geometry} % for adjusting page geometry
\usepackage{titlesec} % for customizing section titles
\usepackage{tabularx} % for making tables with fixed width columns
\usepackage{array} % tabularx requires this
\usepackage[dvipsnames]{xcolor} % for coloring text
\definecolor{primaryColor}{RGB}{0, 0, 0} % define primary color
\usepackage{enumitem} % for customizing lists
\usepackage{fontawesome5} % for using icons
\usepackage{amsmath} % for math
\usepackage[
    pdftitle={John Doe's CV},
    pdfauthor={John Doe},
    pdfcreator={LaTeX with RenderCV},
    colorlinks=true,
    urlcolor=primaryColor
]{hyperref} % for links, metadata and bookmarks
\usepackage[pscoord]{eso-pic} % for floating text on the page
\usepackage{calc} % for calculating lengths
\usepackage{bookmark} % for bookmarks

\usepackage{changepage} % for one column entries (adjustwidth environment)
\usepackage{paracol} % for two and three column entries
\usepackage{ifthen} % for conditional statements
\usepackage{needspace} % for avoiding page brake right after the section title
\usepackage{iftex} % check if engine is pdflatex, xetex or luatex

% Ensure that generate pdf is machine readable/ATS parsable:
\ifPDFTeX
    \input{glyphtounicode}
    \pdfgentounicode=1
    \usepackage[T1]{fontenc}
    \usepackage[utf8]{inputenc}
    \usepackage{lmodern}
\fi

\usepackage{charter}

% Some settings:
\raggedright
\AtBeginEnvironment{adjustwidth}{\partopsep0pt} % remove space before adjustwidth environment
\pagestyle{empty} % no header or footer
\setcounter{secnumdepth}{0} % no section numbering
\setlength{\parindent}{0pt} % no indentation
\setlength{\topskip}{0pt} % no top skip
\setlength{\columnsep}{0.15cm} % set column seperation
\pagenumbering{gobble} % no page numbering

\titleformat{\section}{\needspace{4\baselineskip}\bfseries\large}{}{0pt}{}[\vspace{1pt}\titlerule]

\titlespacing{\section}{
    % left space:
    -1pt
}{
    % top space:
    0.3 cm
}{
    % bottom space:
    0.2 cm
} % section title spacing

\renewcommand\labelitemi{$\vcenter{\hbox{\small$\bullet$}}$} % custom bullet points
\newenvironment{highlights}{
    \begin{itemize}[
        topsep=0.10 cm,
        parsep=0.10 cm,
        partopsep=0pt,
        itemsep=0pt,
        leftmargin=0 cm + 10pt
    ]
}{
    \end{itemize}
} % new environment for highlights


\newenvironment{highlightsforbulletentries}{
    \begin{itemize}[
        topsep=0.10 cm,
        parsep=0.10 cm,
        partopsep=0pt,
        itemsep=0pt,
        leftmargin=10pt
    ]
}{
    \end{itemize}
} % new environment for highlights for bullet entries

\newenvironment{onecolentry}{
    \begin{adjustwidth}{
        0 cm + 0.00001 cm
    }{
        0 cm + 0.00001 cm
    }
}{
    \end{adjustwidth}
} % new environment for one column entries

\newenvironment{twocolentry}[2][]{
    \onecolentry
    \def\secondColumn{#2}
    \setcolumnwidth{\fill, 4.5 cm}
    \begin{paracol}{2}
}{
    \switchcolumn \raggedleft \secondColumn
    \end{paracol}
    \endonecolentry
} % new environment for two column entries

\newenvironment{threecolentry}[3][]{
    \onecolentry
    \def\thirdColumn{#3}
    \setcolumnwidth{, \fill, 4.5 cm}
    \begin{paracol}{3}
    {\raggedright #2} \switchcolumn
}{
    \switchcolumn \raggedleft \thirdColumn
    \end{paracol}
    \endonecolentry
} % new environment for three column entries

\newenvironment{header}{
    \setlength{\topsep}{0pt}\par\kern\topsep\centering\linespread{1.5}
}{
    \par\kern\topsep
} % new environment for the header

% save the original href command in a new command:
\let\hrefWithoutArrow\href

\begin{document}
    \newcommand{\AND}{\unskip
        \cleaders\copy\ANDbox\hskip\wd\ANDbox
        \ignorespaces
    }
    \newsavebox\ANDbox
    \sbox\ANDbox{$|$}

    \begin{header}
        \fontsize{25 pt}{25 pt}\selectfont Sayak Roy\\
        \fontsize{13pt}{15pt}\selectfont AI Infrastructure \& Automation Engineer
        \vspace{0.2 cm}

        \normalsize
        \mbox{Chennai}%
        \kern 5.0 pt%
        \AND%
        \kern 5.0 pt%
        \mbox{\hrefWithoutArrow{mailto:sayakr05@gmail.com}{Gmail}}%
        \kern 5.0 pt%
        \AND%
        \kern 5.0 pt%
        \mbox{\hrefWithoutArrow{tel:91 6033073031}{6033073031}}%
        \kern 5.0 pt%
        \AND%
        \kern 5.0 pt%
        \mbox{\hrefWithoutArrow{https://sayakroyportfolio.netlify.app/}{Portfolio}}%
        \kern 5.0 pt%
        \AND%
        \kern 5.0 pt%
        \mbox{\hrefWithoutArrow{https://www.linkedin.com/in/sayak--roy}{linkedin.com}}%
        \kern 5.0 pt%
        \AND%
        \kern 5.0 pt%
        \mbox{\hrefWithoutArrow{https://github.com/sayakr428}{github.com}}%
    \end{header}

    \vspace{5 pt - 0.3 cm}

    \section{Professional Summary}
\begin{onecolentry}
At \textbf{Tata Communications}, I work as an \textbf{AI Infrastructure \& Automation Engineer} with 3+ years of experience across \textbf{cloud}, \textbf{DevOps}, and \textbf{network automation}. I have been developing \textbf{AI-integrated infrastructure systems} using \textbf{Terraform}, \textbf{Ansible}, and applying my \textbf{full-stack engineering} knowledge. I construct \textbf{automation systems} for \textbf{AWS}, \textbf{Docker}, and \textbf{Kubernetes}, and hold practical understanding of \textbf{BGP}, \textbf{SDN}, and \textbf{telemetry systems}. I am passionate about \textbf{building automation platforms} that integrate \textbf{DevOps with AI tools} to create scalable and intelligent systems.
\end{onecolentry}
        \vspace{0.2 cm}

 \section{Skills}
\begin{onecolentry}
    \begin{highlightsforbulletentries}
        \item \textbf{Cloud \& Infrastructure} – AWS (EC2, Lambda, S3, CloudWatch, API Gateway, EKS, ECS), Azure, Terraform, Ansible, AWS CDK, CloudFormation, NGINX, Docker, Kubernetes.
        \item \textbf{Programming \& Frameworks} – \textbf{Python}, \textbf{JavaScript}, HTML, MATLAB, \textbf{Streamlit}, \textbf{FAST API}, \textbf{Flask}, \textbf{Selenium}, Power Automate.
        \item \textbf{Operating Systems} – Linux (Ubuntu, Debian), Windows, MAC OS.
        \item \textbf{DevOps \& CI/CD} – GitHub Actions, Jenkins, GitLab CI/CD, Cloud-native deployment.
        \item \textbf{Machine Learning \& AI} – Time-Series Models (\textbf{ARIMA}, \textbf{SARIMAX}, \textbf{Prophet}, \textbf{XGBoost}, \textbf{LSTM}, \textbf{TimeGPT}), Ensemble Learning, Anomaly Detection (\textbf{Isolation Forest, LOF, One-Class SVM}), Clustering (\textbf{K-Means}, \textbf{DBSCAN}), \textbf{CNN}, \textbf{RNN}.
        \item \textbf{Full-Stack Development} – React.js, Tailwind CSS, Express.js, Flask, RESTful API Design, JWT Auth, Environment Management, PostgreSQL, MongoDB, SQL.
        \item \textbf{Data Engineering \& Automation} – \textbf{Pandas}, \textbf{NumPy}, \textbf{OpenPyXL}, \textbf{XlsxWriter}, \textbf{python-pptx}, \textbf{ETL Pipelines}, Excel/PPT Automation, MIS Reporting, API Integrations.
        \item \textbf{Monitoring \& Visualization} – Grafana, Prometheus, Kentik, Power BI, Chronograf, Telegraf, CloudWatch.
        \item \textbf{Tools \& Platforms} – \textbf{Git}, \textbf{GitHub}, Power Automate, SharePoint, \textbf{n8n}, \textbf{LangChain}, \textbf{LangGraph}, \textbf{ChromaDB}, \textbf{FAISS}, \textbf{Pinecone}, MS Forms.
    \end{highlightsforbulletentries}
\end{onecolentry}

        \section{Education}
\begin{twocolentry}{}{}
    \textbf{SRM University}, B.Tech in Computer Science (Specialization: Cyber Security)
\end{twocolentry}
        \vspace{0.2 cm}

    \section{Experience}
\begin{twocolentry}{
    December 2025 – Present
}
    \textbf{Deployment \& Maintenance Software Engineer}, DERQ -- Dubai, UAE (Remote)
\end{twocolentry}
\vspace{0.10 cm}
\begin{onecolentry}
    \begin{highlights}
        \item Deploy, configure, and maintain AI-powered Intelligent Transportation Systems (ITS) across distributed edge environments using \textbf{Linux}, \textbf{Docker}, and \textbf{Python}.
        \item Troubleshoot production issues involving \textbf{Edge Processors}, \textbf{RTSP/IP cameras}, \textbf{VPN}, \textbf{TCP/IP}, \textbf{SSH}, and application services through log analysis and root cause identification.
        \item Develop Python automation tools for deployment, system validation, monitoring, and operational maintenance, improving deployment efficiency and reducing manual effort.
        \item Collaborate with software, AI, and infrastructure teams to perform remote deployments, software upgrades, configuration management, and production support across global customer sites.
    \end{highlights}
\end{onecolentry}
\vspace{0.20cm}
\begin{twocolentry}{
    April 2023 – December 2025
}
    \textbf{AI Infrastructure \& Automation Engineer}, Tata Communications -- Chennai, Tamil Nadu
\end{twocolentry}
\vspace{0.10 cm}
\begin{onecolentry}
    \begin{highlights}
        \item Automated router configurations using \textbf{OpenDaylight SDN Controller} with \textbf{NETCONF/RESTCONF APIs}, reducing manual provisioning effort.
        \item Designed YAML-based \textbf{OpenConfig templates} for multi-vendor routing protocols (BGP, OSPF, IS-IS, EIGRP).
        \item Built \textbf{RESTful APIs} to fetch live routing data and push configs, improving operational efficiency and reliability by almost 30\%.
        \item Automated \textbf{IRR-driven Route Paths and ASN List/AS Set} using \textbf{bgpq3} and custom Python scripts to fetch, parse prefix lists and as-path sets from IRR sources, reduced manual effort by 80\%.
        \item Built and deployed advanced forecasting systems using \textbf{ARIMA, SARIMAX, Prophet, XGBoost, LSTM, TimeGPT}, enabling accurate traffic prediction, capacity planning, and anomaly detection across Ethernet \& ILL backbones.
        \item Designed production-grade ML pipelines with dynamic hyperparameter tuning, feature engineering, \textbf{Boosting Ensembles} (Prophet + XGBoost), and a \textbf{3-layer Stacking Ensemble} (\textbf{XGBoost + LSTM + Meta-Model}), successfully deployed for real-time inference.
        \item Developed anomaly detection engines using \textbf{Z-Score, IQR, Isolation Forest, LOF, One-Class SVM} and clustering models including \textbf{K-Means} and \textbf{DBSCAN} for customer segmentation and traffic behavior analysis.
    \end{highlights}
\end{onecolentry}
\vspace{0.2 cm}
\begin{twocolentry}{
    March 2022 – Jan 2023
}
    \textbf{Analyst Intern}, Accenture -- Bangalore, Karnataka
\end{twocolentry}
\vspace{0.10 cm}
\begin{onecolentry}
    \begin{highlights}
        \item Developed \textbf{dynamic dashboards} and \textbf{data-driven UI} using \textbf{React.js}, \textbf{Tailwind CSS}, and \textbf{Redux-Saga}, enabling seamless visualization of campaign performance and engagement metrics.
        \item Utilized \textbf{Python}, \textbf{MongoDB}, and \textbf{SQL} for querying, preprocessing, and implementing \textbf{dynamic data filtering} powering interactive visual elements.
        \item Worked with \textbf{cross-functional teams} to define KPIs, optimize \textbf{API performance}, and design efficient \textbf{JSON payload structures} for streamlined communication.
    \end{highlights}
\end{onecolentry}

    \section{Certifications}
\begin{onecolentry}
    \begin{highlights}
        \item Google IT Automation with Python Professional Certificate.
        \item Building Modern Python Applications on AWS.
        \item Designing and Implementing Azure AI Solutions.
    \end{highlights}
\end{onecolentry}

 \section{Projects}
\begin{onecolentry}
    \textbf{\href{https://github.com/sayakr428/leo-kpi-monitoring}{LEO Satellite KPI Monitoring using mtr, iPerf3 \& InfluxDB}}
    \begin{highlights}
        \item Deployed \textbf{Raspberry Pi} devices as low‑cost, on‑site network probes at multiple test locations.
        \item Integrated with \textbf{LEO satellite} terminals to capture end‑to‑end performance data.
        \item Automated KPI collection using tools like \textbf{iperf3} and \textbf{mtr}, scheduled via \textbf{cron‑jobs}.
        \item Built \textbf{Python‑scripts} to process metrics and send them to a central \textbf{InfluxDB} time‑series database.
        \item Configured \textbf{Telegraf} for secure, periodic data push from probes to the monitoring server.
        \item Developed real‑time dashboards in \textbf{Chronograf} for visualization and alerting.
        \item Enabled historical performance tracking and geographical comparison of satellite link quality.
    \end{highlights}
\end{onecolentry}
\vspace{0.20 cm}
\begin{onecolentry}
    \textbf{\href{https://github.com/sayakr428/Implemented-BGP-monitoring-Protocol-to-get-graphical-representation-of-the-BGP-data-s}{BMP to get graphical representation of the BGP data’s}}
    \begin{highlights}
        \item Designed and deployed a containerized monitoring framework to collect, process, and visualize BGP routing updates using the \textbf{BGP Monitoring Protocol}.
        \item Deployed \textbf{OpenBMP} collector and supporting services in \textbf{Docker} containers for portability and ease of deployment.
        \item Integrated with routers to receive BMP feeds containing real‑time BGP routing table snapshots and updates.
        \item Configured \textbf{Zookeeper} \& \textbf{Kafka} within Docker for scalable BMP message streaming.
        \item Stored and indexed routing data in \textbf{PostgreSQL} for historical analysis.
        \item Built \textbf{Grafana} dashboard for monitoring BGP peer health, routing stability, and prefix reachability.
        \item Automated alerts for route flaps, \textbf{suspicious prefix} announcements, and dropouts from the peer session.
    \end{highlights}
\end{onecolentry}
\vspace{0.20cm}
\begin{onecolentry}
    \textbf{\href{https://github.com/sayakr428/cloud-butler-infra}{AI Chatbot \& n8n Automation Pipeline}}
    \begin{highlights}
        \item Built a chat-driven DevOps automation system using \textbf{n8n}, \textbf{OpenAI}, and \textbf{Terraform Cloud}, enabling users to provision AWS infrastructure through natural language commands.
        \item Designed AI workflows that convert user prompts into validated \textbf{Terraform templates} for EC2, S3, RDS, VPC, IAM, and Security Groups, ensuring fully automated, error-free provisioning.
        \item Integrated \textbf{remote Terraform workspaces} with dynamic variable injection, automated runs, and real-time status updates returned directly to the chatbot interface.
        \item Delivered an end-to-end \textbf{Conversational Infrastructure-as-Code (CIaC)} pipeline that eliminated manual configuration and demonstrated advanced cloud automation and AI orchestration capabilities.
    \end{highlights}
\end{onecolentry}

    \section{ Competitions And Awards}
        \begin{onecolentry}
    \textbf{Extra Mile Award} \hfill \textit{Aug 2024 | Tata Communications}
    \begin{highlights}
        \item Received the “\textbf{Extra Mile Award}” for designing and developing an automation framework that dynamically generated \textbf{JSON configuration files} for diverse \textbf{BGP session setups}.
        \item Integrated the solution with Tata Communications’ private \textbf{IZO Cloud platform}, enabling seamless deployment via \textbf{Ansible playbooks}.
    \end{highlights}
\end{onecolentry}

\end{document}
